%!TEX TS-program = xelatex
%!TEX encoding = UTF-8 Unicode
% Awesome CV LaTeX Template for CV/Resume
%
% This template has been downloaded from:
% https://github.com/posquit0/Awesome-CV
%
% Author:
% Claud D. Park <posquit0.bj@gmail.com>
% http://www.posquit0.com
%
% Template license:
% CC BY-SA 4.0 (https://creativecommons.org/licenses/by-sa/4.0/)
%


%-------------------------------------------------------------------------------
% CONFIGURATIONS
%-------------------------------------------------------------------------------
% A4 paper size by default, use 'letterpaper' for US letter
\documentclass[11pt, a4paper]{awesome-cv}




% Configure page margins with geometry
\usepackage{needspace}
\geometry{left=1.4cm, top=.8cm, right=1.4cm, bottom=1.5cm, footskip=.5cm}

% Specify the location of the included fonts
\fontdir[fonts/]

% Color for highlights
% Awesome Colors: awesome-emerald, awesome-skyblue, awesome-red, awesome-pink, awesome-orange
%                 awesome-nephritis, awesome-concrete, awesome-darknight
% \colorlet{awesome}{awesome-red}
% Uncomment if you would like to specify your own color
% \definecolor{awesome}{HTML}{1E4E8C}
\definecolor{awesome}{HTML}{2563EB}

% Korean: switch to Noto Sans CJK KR for Hangul automatically (XeTeX interchar classes)
\newfontfamily\hangulfont{Noto Sans CJK KR}[Scale=0.94]
\newfontfamily\hangulfontbold{Noto Sans CJK KR Medium}[Scale=0.94]
\XeTeXinterchartokenstate=1
\newXeTeXintercharclass\HangulClass
\count255="AC00 \loop\ifnum\count255<"D7A4 \XeTeXcharclass\count255=\HangulClass \advance\count255 by 1 \repeat
\count255="3131 \loop\ifnum\count255<"318F \XeTeXcharclass\count255=\HangulClass \advance\count255 by 1 \repeat
\XeTeXinterchartoks 0 \HangulClass = {\begingroup\hangulfont}
\XeTeXinterchartoks 4095 \HangulClass = {\begingroup\hangulfont}
\XeTeXinterchartoks \HangulClass 0 = {\endgroup}
\XeTeXinterchartoks \HangulClass 4095 = {\endgroup}

% Colors for text
% Uncomment if you would like to specify your own color
% \definecolor{darktext}{HTML}{414141}
% \definecolor{text}{HTML}{333333}
% \definecolor{graytext}{HTML}{5D5D5D}
% \definecolor{lighttext}{HTML}{999999}

% Set false if you don't want to highlight section with awesome color
\setbool{acvSectionColorHighlight}{true}
\renewcommand{\acvSectionTopSkip}{0mm}
\renewcommand{\acvSectionContentTopSkip}{1.5mm}

% If you would like to change the social information separator from a pipe (|) to something else
\renewcommand{\acvHeaderSocialSep}{\quad\textbar\quad}


%-------------------------------------------------------------------------------
%	PERSONAL INFORMATION
%	Comment any of the lines below if they are not required
%-------------------------------------------------------------------------------
% Available options: circle|rectangle,edge/noedge,left/right
% \photo{./examples/profile.png}
\name{노}{장현}
\renewcommand{\acvHeaderNameDelim}{}
\position{AI 연구자}
% \address{42-8, Bangbae-ro 15-gil, Seocho-gu, Seoul, 00681, Rep. of KOREA}

\email{watercar99@naver.com}
\homepage{janghyunnoh.github.io}
\github{janghyunnoh}
\linkedin{janghyun-noh-411587195}
% \gitlab{gitlab-id}
% \stackoverflow{SO-id}{SO-name}
% \twitter{@twit}
% \skype{skype-id}
% \reddit{reddit-id}
% \medium{madium-id}
\googlescholar{NPIG8VoAAAAJ}{Google Scholar}
%% \firstname and \lastname will be used
% \googlescholar{googlescholar-id}{}
% \extrainfo{extra informations}

% \quote{``Be the change that you want to see in the world."}


%-------------------------------------------------------------------------------
\let\origcvsection\cvsection
\renewcommand{\cvsection}[1]{\needspace{7\baselineskip}\origcvsection{#1}}
\let\origcvsubsection\cvsubsection
\renewcommand{\cvsubsection}[1]{\needspace{5\baselineskip}\origcvsubsection{#1}}
\begin{document}

% Print the header with above personal informations
% Give optional argument to change alignment(C: center, L: left, R: right)
\makecvheader\vspace{4mm}

% Print the footer with 3 arguments(<left>, <center>, <right>)
% Leave any of these blank if they are not needed
\makecvfooter
  {\the\year 년 \the\month 월 \the\day 일}
  {노장현~~~·~~~이력서}
  {\thepage}


%-------------------------------------------------------------------------------
%	CV/RESUME CONTENT
%	Each section is imported separately, open each file in turn to modify content
%-------------------------------------------------------------------------------


%-------------------------------------------------------------------------------
%	SECTION TITLE
%-------------------------------------------------------------------------------
\cvsection{학력}


%-------------------------------------------------------------------------------
%	CONTENT
%-------------------------------------------------------------------------------
\begin{cventries}
\cventry
    {인공지능학 석사}
    {명지대학교}
    {서울}
    {2024.09 -- 2026.08}
    {
      \begin{cvitems}
        \item {인공지능 및 데이터 분석(AIDA) 연구실에서 오민식 교수님 지도 아래 연구 수행}
        \item {서울대학교병원, 질병관리청, 경북대학교, (주)아이겐드럭 등 외부 기관과 공동 연구 수행}
        \item {딥러닝, 머신러닝, 데이터마이닝, 프로그래밍 등 학부 과목 조교 (2024년 2학기 -- 2026년 1학기)}
      \end{cvitems}
    }
\vspace{-1mm}
\cventry
    {융합소프트웨어학부 데이터테크놀로지 전공 학사}
    {명지대학교}
    {서울}
    {2018.03 -- 2024.08}
    {
      \begin{cvitems}
        \item {인공지능 및 데이터 분석(AIDA) 연구실 학부연구생으로 오민식 교수님 지도 아래 연구 수행}
        \item {질병관리청과 고지혈증 환자의 고혈압 진행 관련 임상·유전적 요인 공동 연구}
      \end{cvitems}
    }
\vspace{-1mm}
\cventry
    {재학}
    {레드스쿨 (대안 중고등학교)}
    {금산}
    {2012.03 -- 2018.01}
    {
      \begin{cvitems}
        \item {2년 연속 성적 우수 장학금 수혜}
      \end{cvitems}
    }
\end{cventries}

%-------------------------------------------------------------------------------
%	SECTION TITLE
%-------------------------------------------------------------------------------
\vspace{-3.5mm}\cvsection{연구 경력}


%-------------------------------------------------------------------------------
%	CONTENT
%-------------------------------------------------------------------------------
\begin{cventries}
\cventry
    {학부연구생 / 석사과정 연구원} % Job title
    {명지대학교 인공지능 및 데이터 분석(AIDA) 연구실} % Organization
    {서울} % Location
    {2022.12 -- 2026.08} % Date(s)
    {
      \begin{cvitems} % Description(s) of tasks/responsibilities     
        \item {단일세포 RNA-seq 데이터의 세포 주석 없는 표현형 예측 프레임워크 scCap 개발 및 코드·Docker 이미지 공개 (Briefings in Bioinformatics 1저자)}
        \item {scRNA-seq 기반 MSI/MSS 표현형 예측 모델 개발 및 AI 기반 타겟 발굴 연구 (아이겐드럭)}
        \item {TNBC Xenium 공간전사체 데이터 기반 세포 niche 및 종양 미세환경(TME) 분석 (서울대학교병원)}
        \item {HCC GeoMx 공간전사체 기반 간경변·비간경변 조직 비교 및 면역조절 타겟 발굴}
        \item {CRISPR-Cas9 유도 대규모 DNA 결실 분석을 위한 Biomni 기반 AI 에이전트 환경 구축}
        \item {MIMIC-IV 임상 데이터와 MIMIC-CXR 흉부 X-ray 기반 코로나19 중증도 예측 멀티모달 모델 개발 (질병관리청)}
        \item {사과 전장 유전체 데이터를 위한 nf-core/sarek 변이 분석 파이프라인 구축 (경북대학교)}
        \item {KoGES 코호트 데이터의 수치 기반 라벨링 설계 및 고혈압 진행 예측 모델 개발 (질병관리청)}
      \end{cvitems}
    }
\end{cventries}


%-------------------------------------------------------------------------------
%	SECTION TITLE
%-------------------------------------------------------------------------------
\vspace{-3.5mm}\cvsection{논문 및 학회 발표}

%-------------------------------------------------------------------------------
%	CONTENT
%-------------------------------------------------------------------------------

%-------------------------------------------------------------------------------
%	JOURNAL ARTICLES
%-------------------------------------------------------------------------------
\cvsubsection{저널 논문}

\begin{cventries}
\cventry
    {게재}
    {Annotation-free phenotype prediction using knowledge-augmented clustering \newline from single-cell RNA sequencing data}
    {Briefings in Bioinformatics (BIB)}
    {2026}
    {
      \textbf{J. Noh}, Y. Shin, M. Kim, and M. Oh
    }
\vspace{3.5mm}
\cventry
    {게재}
    {Ferroptosis-based therapeutic strategy targeting iron metabolism and SLC7A11 overcomes oxaliplatin resistance in patient-derived pancreatic cancer organoids}
    {Journal of Experimental \& Clinical Cancer Research (JECCR)}
    {2026}
    {
      J. Choi, K. Lee, M. Oh, \textbf{J. Noh}, Y. Kim, E. Hwang,
      I. Cho, W. Paik, J. Ryu, J. Ku, and S. Lee
    }
\end{cventries}

%-------------------------------------------------------------------------------
%	CONFERENCE PRESENTATIONS
%-------------------------------------------------------------------------------
\cvsubsection{학회 발표}

\begin{cventries}
\cventry
    {구두 발표}
    {scRNA-seq 기반 표현형 예측을 위한 사전 학습 임베딩을 활용한 환자 단위 메타셀 구성}
    {한국정보과학회 \newline 한국컴퓨터종합학술대회}
    {2026}
    {
      Y. Shin, \textbf{J. Noh}, and M. Oh
    }
\vspace{3.5mm}
\cventry
    {포스터 발표 (우수 포스터상)}
    {scRNA-seq 데이터 축소 기법을 이용한 환자 표현형 예측의 연산 경량화 분석}
    {한국정보과학회 \newline 한국소프트웨어종합학술대회}
    {2025}
    {
      Y. Shin, \textbf{J. Noh}, and M. Oh
    }
\vspace{3.5mm}
\cventry
    {포스터 발표}
    {An annotation-free framework for phenotype prediction using hierarchical multiple instance learning with single-cell RNA-sequencing data}
    {International Conference on Genome Informatics ISCB-Asia}
    {2025}
    {
      \textbf{J. Noh}, Y. Shin, M. Kim, and M. Oh
    }
\vspace{3.5mm}
\cventry
    {포스터 발표 (우수 포스터상)}
    {Patient Phenotype Prediction from Single-cell RNA-seq Data through Annotation-free Multiple Instance Learning}
    {한국생명정보학회 \newline BIOINFO 2025}
    {2025}
    {
      \textbf{J. Noh}, Y. Shin, M. Kim, and M. Oh
    }
\vspace{3.5mm}
\cventry
    {구두 발표 (우수발표논문상)}
    {AMIL: Automated Cell Grouping Framework for Phenotype Prediction with Multiple Instance Learning}
    {한국정보과학회 \newline 한국컴퓨터종합학술대회}
    {2025}
    {
      \textbf{J. Noh}, M. Kim, and M. Oh
    }
\vspace{3.5mm}
\cventry
    {포스터 발표}
    {MIMIC-III 데이터를 이용한 급성호흡곤란증후군 환자의 병원 내 사망 예측: 다양한 결측 데이터 처리 전략을 적용한 머신러닝 모델의 비교 연구}
    {한국정보과학회 \newline 한국소프트웨어종합학술대회}
    {2023}
    {
      \textbf{J. Noh}, G. Choi, and M. Oh
    }
\end{cventries}


%-------------------------------------------------------------------------------
%	SECTION TITLE
%-------------------------------------------------------------------------------
\vspace{1mm}\cvsection{수상 및 장학}

%-------------------------------------------------------------------------------
%	SUBSECTION TITLE
%-------------------------------------------------------------------------------
\cvsubsection{수상}

%-------------------------------------------------------------------------------
%	CONTENT
%-------------------------------------------------------------------------------
\begin{cvhonors}

%---------------------------------------------------------
  \cvhonor
    {우수연구상}
    {명지대학교}
    {서울}
    {2026}
%---------------------------------------------------------
  \cvhonor
    {우수 포스터상}
    {한국정보과학회 한국소프트웨어종합학술대회}
    {서울}
    {2025}
%---------------------------------------------------------
  \cvhonor
    {우수 포스터상}
    {한국생명정보학회 BIOINFO 2025}
    {서울}
    {2025}
%---------------------------------------------------------
  \cvhonor
    {우수발표논문상}
    {한국정보과학회 한국컴퓨터종합학술대회}
    {서울}
    {2025}

%---------------------------------------------------------
  \cvhonor
    {기관장 표창장}
    {대한무역투자진흥공사 (KOTRA)}
    {서울}
    {2022}

%---------------------------------------------------------
\end{cvhonors}

%-------------------------------------------------------------------------------
%	SUBSECTION TITLE
%-------------------------------------------------------------------------------
\cvsubsection{장학}

%-------------------------------------------------------------------------------
%	CONTENT
%-------------------------------------------------------------------------------


%---------------------------------------------------------

\begin{cvhonors}

\cvhonor
    {TriBIG 프로젝트 장학금}
    {대한무역투자진흥공사 (KOTRA)}
    {서울}
    {2022}
%---------------------------------------------------------

  \cvhonor
    {세계화 장학금}
    {명지대학교}
    {서울}
    {2019}

%---------------------------------------------------------
  \cvhonor
    {신입생 성적 우수 장학금 3종}
    {명지대학교}
    {서울}
    {2018}

\cvhonor
    {성적 우수 장학금}
    {레드스쿨}
    {금산}
    {2016}

\cvhonor
    {성적 우수 장학금}
    {레드스쿨}
    {금산}
    {2015}

%---------------------------------------------------------

\end{cvhonors}

%-------------------------------------------------------------------------------
%	SECTION TITLE
%-------------------------------------------------------------------------------
\needspace{5\baselineskip}
\vspace{1mm}\cvsection{프로젝트}

%-------------------------------------------------------------------------------
%	CONTENT
%-------------------------------------------------------------------------------
\vspace{\acvSectionContentTopSkip}\vspace{-2.0mm}
\setlength\tabcolsep{1ex}\setlength{\extrarowheight}{0pt}
\noindent\begin{tabularx}{\textwidth}{l>{\raggedright\arraybackslash}X}
  \skilltypestyle{\makebox[1.9cm][l]{2024 -- 2026}} & \leavevmode\skillsetstyle{AI 기반 타겟 발굴 플랫폼 확립을 통한 만성 염증성 질환 조절 CAR-Treg 세포 치료제 개발} \\
  \skilltypestyle{\makebox[1.9cm][l]{2024 -- 2026}} & \leavevmode\skillsetstyle{기후 변화 대응을 위한 사과꽃 휴면타파 및 꽃 발달 시기의 유전적 생리학적 메커니즘 연구} \\
  \skilltypestyle{\makebox[1.9cm][l]{2024 -- 2026}} & \leavevmode\skillsetstyle{공간 전사체 기반 타겟 발굴과 오가노이드 검증을 통한 난치성 고형암 신규 치료제 개발} \\
  \skilltypestyle{\makebox[1.9cm][l]{2025 -- 2026}} & \leavevmode\skillsetstyle{CRISPR-Cas9 유도 대규모 DNA 결실 검출·정량화 웹 분석 플랫폼 개발} \\
  \skilltypestyle{\makebox[1.9cm][l]{2025 -- 2026}} & \leavevmode\skillsetstyle{HCC GeoMx 공간전사체 기반 간경변·비간경변 조직의 병리적 차이 규명} \\
  \skilltypestyle{\makebox[1.9cm][l]{2025}} & \leavevmode\skillsetstyle{멀티모달 임상·영상 데이터를 활용한 코로나19 환자 중증도 예측} \\
  \skilltypestyle{\makebox[1.9cm][l]{2023 -- 2024}} & \leavevmode\skillsetstyle{KoGES 코호트 기반 고지혈증 환자의 고혈압 진행 위험군 규명} \\
  \skilltypestyle{\makebox[1.9cm][l]{2023}} & \leavevmode\skillsetstyle{그놈을 잡아라 (CatchYou): 진술 기반 몽타주 생성 서비스} \\
  \skilltypestyle{\makebox[1.9cm][l]{2022}} & \leavevmode\skillsetstyle{글로벌 공급망 위기 대응을 위한 세관 고시 분석 및 모니터링 시스템 개발} \\
\end{tabularx}\par

%-------------------------------------------------------------------------------
\vspace{1mm}\cvsection{기술 역량}

%-------------------------------------------------------------------------------
%	CONTENT
%-------------------------------------------------------------------------------
\begin{cvskills}

%---------------------------------------------------------
  \cvskill
    {언어}
    {Python, R}

%---------------------------------------------------------
  \cvskill
    {프레임워크}
    {PyTorch, TensorFlow, Scikit-learn, XGBoost, Scanpy, FastAPI}

%---------------------------------------------------------
  \cvskill
    {도구}
    {Linux, Git, Docker, Nextflow}

%---------------------------------------------------------
  \cvskill
    {연구 분야}
    {머신러닝, 딥러닝, 멀티모달 학습, 의료 AI, 바이오인포매틱스}

%---------------------------------------------------------
  \cvskill
    {외국어}
    {한국어(모국어), 영어(업무 능통; OPIc IH 2026, TOEIC 920점 2023)}

%---------------------------------------------------------
  \cvskill
    {자격증}
    {정보처리기사 (2023), SQL 개발자 SQLD (2023)}

%---------------------------------------------------------
\end{cvskills}

%-------------------------------------------------------------------------------
%	SECTION TITLE
%-------------------------------------------------------------------------------
\vspace{1mm}\cvsection{대외활동}

%-------------------------------------------------------------------------------
%	SUBSECTION TITLE
%-------------------------------------------------------------------------------
\cvsubsection{국내}

%-------------------------------------------------------------------------------
%	CONTENT
%-------------------------------------------------------------------------------
\begin{cventries}
\cventry
    {AI 팀원}
    {Google Developer Student Clubs, 명지대학교}
    {서울}
    {2022.10 -- 2023.12}
    {
      \begin{cvitems}
        \item {머신러닝·딥러닝 스터디 활동}
        \item {AI 프로젝트 협업 및 AI 경진대회 참가}
      \end{cvitems}
    }
\vspace{-1mm}
\cventry
    {프로젝트 팀원}
    {빅리더 AI 아카데미}
    {서울}
    {2022.07 -- 2022.10}
    {
      \begin{cvitems}
        \item {KOTRA 빅데이터팀과 세관 고시 분석 및 모니터링 시스템 개발}
        \item {TriBIG 프로젝트 우수 성과로 장학금 수혜}
      \end{cvitems}
    }
\vspace{-1mm}
\cventry
    {참가자}
    {데이터 청년 캠퍼스, 동국대학교}
    {서울}
    {2022.07 -- 2022.08}
    {
      \begin{cvitems}
        \item {한국데이터산업진흥원 주관 데이터 사이언스 기반 지능형 소프트웨어 과정 수료}
      \end{cvitems}
    }
\end{cventries}

%-------------------------------------------------------------------------------
%	SUBSECTION TITLE
%-------------------------------------------------------------------------------
\cvsubsection{국외}

%-------------------------------------------------------------------------------
%	CONTENT
%-------------------------------------------------------------------------------
\begin{cventries}
\cventry
    {순례자}
    {산티아고 성지순례 (프랑스길)}
    {스페인}
    {2026.09 -- 2026.10}
    {}
\vspace{3.5mm}
\cventry
    {봉사자}
    {필리핀 의료선교 (GMMA)}
    {필리핀 마닐라}
    {2024.01}
    {
      \begin{cvitems}
        \item {국제 의료진과 함께 의료 봉사 활동 참여}
      \end{cvitems}
    }
\vspace{-1mm}
\cventry
    {봉사자}
    {아이슬란드 국제 워크캠프}
    {아이슬란드 레이캬비크}
    {2019.12 -- 2020.01}
    {
      \begin{cvitems}
        \item {환경 캠페인 및 국제 봉사자들과의 문화 교류 활동}
      \end{cvitems}
    }
\vspace{-1mm}
\cventry
    {봉사자}
    {베트남 단기선교}
    {베트남 하노이}
    {2018.08}
    {
      \begin{cvitems}
        \item {현지 주민·대학생과의 문화 교류 및 지역 봉사 활동}
      \end{cvitems}
    }
\end{cventries}

\end{document}
